\section{What integer corrections can attain}\label{sec:corrections}

The remaining question is how much the arithmetic congruences can
increase the determinant. We examine two ways to control that increase:
describe the least attainable nonzero value under the rationality
assumption, or construct corrections with a sufficiently small real
effect. The first approach leads to a row-lattice covolume; the second
leads to rounding in the correction's output space. We then test
polynomial kernel inverses as a way to bound the covolume, and explain
why successful finite searches do not settle the uniform question.

\subsection{An arithmetic progression of determinant values}

The correction matrix has many entries, but its determinant values
have a simple exact description at a rational target. This lets us
replace a search over large integer matrices by a question about the
step of one arithmetic progression. The reduction does not control
the height of the matrices needed to attain those values.

We begin with a square congruence problem. If $T\in\Z^{b\times b}$
has Smith invariants $d_1\mid\cdots\mid d_b$, including possible
trailing zeros, and $\ell$ is a positive integer, then
\begin{equation}\label{eq:smith-progression}
 \{\det(T+\ell K):K\in\Z^{b\times b}\}
 =\det T+\ell\prod_{i=1}^{b-1}\gcd(\ell,d_i)\Z.
\end{equation}
The empty product for $b=1$ is one. To prove necessity, expand the
determinant by the number of changed entries. At a prime with
$v_p(\ell)=e$, the smallest possible valuation of a nonconstant term
is $e+\sum_{i<b}\min(e,v_p(d_i))$.
For attainment, pass to Smith form and choose the first $b-1$ rows
upper bidiagonal, with diagonal $d_i$ and superdiagonal $\ell$.
Their last-row cofactors are
\[
 c_j=(-1)^{b+j}\ell^{b-j}d_1\cdots d_{j-1},\qquad 1\le j\le b.
\]
Their gcd is $\prod_{i<b}\gcd(\ell,d_i)$. Varying the last row by
$\ell$ times an integer row and using B\'ezout's identity attains
every value in \eqref{eq:smith-progression}.
For example, matrices congruent to $\diag(2,6)$ modulo $4$ have
determinants exactly $12+8\Z$. The matrices
$\left(\begin{smallmatrix}2&4\\4u&6+4v\end{smallmatrix}\right)$
already attain them all.

For an invertible rational $M$ and a rational $b$ by $n$ matrix $W$
of rank $b$, put $R=WM^{-1}$. Suppose corrections are
$U=U_0+mZ$. A basis matrix $H_\Lambda$ for the rational lattice
$mR\Z^n$ gives
\[
 \det(M+UW)=\det M\det H_\Lambda
       \det\bigl(H_\Lambda^{-1}(I_b+RU_0)+K\bigr),
 \qquad K\in\Z^{b\times b}.
\]
Every $K$ is attainable because the columns vary independently.
Clearing a denominator of the last fixed matrix and applying
\eqref{eq:smith-progression} proves that the determinant values form
an arithmetic progression. Equivalently, its positive step is the
rational gcd of the nonconstant coefficients of
$\det(M+U_0W+mZW)$. These coefficients are mixed row minors; finite
differences recover them from integer-input values. This is the
exact-image result of \rec{R16}.

\subsection{The combined row lattice determines the leading rate}

Here is an intrinsic form that also allows a singular initial matrix.
For finitely many rational numbers, their \emph{rational content}
means the positive generator of their integer span. It is computed by
clearing a denominator and taking an ordinary gcd.

\begin{theorem}\label{thm:row-lattice}
Let $M_0\in\Q^{n\times n}$ and $V\in\Q^{b\times n}$, with
$\rank V=b$ and $\rank\left[\begin{smallmatrix}M_0\\V\end{smallmatrix}\right]=n$.
Choose a row basis $H_\Lambda$ of the lattice generated by the stacked
rows, with positive determinant, and put
\[
 \delta=\det H_\Lambda,\qquad
 A=M_0H_\Lambda^{-1},\quad B=VH_\Lambda^{-1},\quad
 a=\det A,\quad c=\gcd_{i,j}B_{ij}.
\]
Then $A,B$ are integer matrices and
\begin{equation}\label{eq:content-progression}
 \{\det(M_0+ZV):Z\in\Z^{n\times b}\}
       =\delta(a+c\Z),\qquad \gcd(a,c)=1.
\end{equation}
Moreover, $\delta$ is the rational content of all maximal minors
of the stack. Both $\delta$ and $c$ are independent of the chosen basis.
\end{theorem}
\begin{proof}
The rows of $[A;B]$ generate $\Z^n$, so this matrix has an integer
left inverse. Cauchy--Binet shows that its maximal minors have gcd
one. Multiplying back by $H_\Lambda$ proves the assertion about
$\delta$.

Some $M_0+Z_*V$ is invertible: a nonzero mixed minor is a nonzero
coefficient of its determinant polynomial, and a nonzero polynomial
cannot vanish at every integer tuple. Shift the initial matrix by
this $Z_*$ and apply the progression result above. Its step, divided
by $\delta$, is the gcd of maximal minors of $[A;B]$ using at least
one row from $B$.

Fix a prime $p$, and put $e=\min v_p(B_{ij})$. If $e=0$, extend a
nonzero row of $B$ modulo $p$ to a basis using rows of $[A;B]$.
This gives a unit mixed minor. If $e>0$, then $A$ is invertible
modulo $p$. All mixed minors have valuation at least $e$, and the
one-row replacements are entries of $B\adj(A)$, up to sign. Since
$\adj(A)$ is a unit matrix over $\Z_p$, one has valuation exactly
$e$. Thus the step is $c$, prime by prime. The same argument shows
that no prime dividing $c$ divides $a$.
\end{proof}

For example, take $M_0=\diag(1/2,1)$ and $V=(5/3,0)$.
The row lattice has basis $\diag(1/6,1)$, so $\delta=1/6$, $a=3$,
$c=10$. The determinants are $(3+10z)/6$, and their least absolute
value is $1/2$. The lattice covolume is not itself always attained.

If zero is unattainable in \eqref{eq:content-progression}, its least
absolute value $\mu$ satisfies
\begin{equation}\label{eq:least-determinant}
 \delta\le\mu\le\delta c/2.
\end{equation}
If $c=1$, zero is attainable and the least \emph{nonzero} absolute
value is $\delta$. We next show that $c$ is subexponential on the
$n^2$ scale for the moment matrices. This leaves $\delta$ as the
quantity whose leading rate matters.

If an integer $D$ clears every entry of $M_0,V$, then $DH_\Lambda$
is integral. Writing $B=cB'$ gives $DV=cB'(DH_\Lambda)$, so
\begin{equation}\label{eq:step-bound}
 c\le D\max_{i,j}|V_{ij}|.
\end{equation}
For the construction of Section~\ref{sec:moments}, fix $x=u/d\in\Q$
in lowest terms and take $M_0=M_N(x)+U_N^0W$, $V=m_NW$. The bound
on the low-moment rank for $65N/2<p\le145N/4$ gives $\rank W=b$
for all sufficiently large $N$, since that interval then contains a
prime. Proposition~\ref{prop:moment-nonvanishing} makes $M_N(x)$
invertible along large prime scales. Thus the rank hypotheses hold on
the sequence used in the proposed contradiction. A common denominator
\[
 D=d\,2^{H+1}((2H+2)!)^2
\]
clears all entries. Indeed,
$Z^0(i,i)=-\sum_{j=1}^i j^{-2}$ and
$Z^0(i,j)=(H_i-H_j)/(i-j)$ off the diagonal. The rational $M$ moments
satisfy the following recurrences, with $c_j=4^{-j}\binom{2j}j$:
\begin{gather*}
 t_{2j}=2/((2j+1)c_j),\quad t_{2j+1}=0,\quad
 v_0=w_0=w_1=0,\quad v_1=2,\\
 v_k=((k-1)v_{k-2}+t_{k-2}+t_{k-1})/k,\qquad
 w_k=((k-2)w_{k-2}+2/(k-1))/(k-1),\\
 M^0(i,0)=v_i,\quad M^0(0,j)=w_j,\qquad
 M^0(i,j)=M^0(i-1,j-1)-t_{i-1}/j.
\end{gather*}
Expanding a recurrence introduces products of distinct integers no
larger than $2H+2$; the second factorial covers the forcing terms.
The polynomial-part projection has denominator dividing $2^H$.
All other row and column coefficients are integers, proving the claim.

The coefficient sum of $F_i$ is at most $2^{L+2h}$, the column filter
has coefficient sum $2^q$, and $|Z^0(i,j)|<2$. Hence
$\max|W_{ij}|\le48\,2^{L+2h+q}$. Also
$m_N\le2^C(H!)^2$. Formula~\eqref{eq:step-bound} gives
\[
 c\le48d\,2^{H+1+C+L+2h+q}((2H+2)!)^2(H!)^2,
 \qquad \log c=O_x(n\log n).
\]
Together with \eqref{eq:least-determinant}, this proves
\begin{equation}\label{eq:content-rate}
 \log\mu=\log\delta+O_x(n\log n)
\end{equation}
whenever zero is unattainable. This is the all-size conclusion of
\rec{R03}. Under a hypothetical rational value $x=\theta$, a sufficiently
strong estimate for this row-lattice covolume would finish the argument.
No such estimate has been established.

\subsection{A prime-by-prime comparison of the index loss}

The covolume can be compared directly with the canonical determinant.
At a rational $x$ where $M=M_N(x)$ is invertible, put
\[
 \Lambda_1=\operatorname{row}_{\Z}[M;W],\quad
 \Lambda_0=\operatorname{row}_{\Z}[M+U_0W;mW],\quad
 I=[\Lambda_1:\operatorname{row}_{\Z}M],\quad
 J=[\Lambda_1:\Lambda_0].
\]
Then $\delta/|\det M|=J/I$. At $p\nmid m$, the local spans of
$\Lambda_0,\Lambda_1$ agree, so the local contribution to $J/I$ is
nonpositive. If $e=v_p(m)>0$ and $u_i$ are the Smith invariants of
$U_0$ modulo $p^e$, a local diagonalization gives
\[
 v_p(J/I)\le\sum_{i=1}^b(e-v_p(u_i))_+.
\]
Indeed, multiplying the affected canonical rows by the indicated
powers of $p$ puts them in $\Lambda_0$; this bounds the loss of index.
For the choices in Theorem~\ref{thm:arithmetic-cost}, only primes
$11N<p\le52N$ contribute positively. The bounds are $10N$ up to
$65N/2$ and $14N$ afterwards. The prime number theorem therefore gives
\[
 \log(J/I)\le(488+o(1))N^2=(61/288+o(1))n^2.
\]
This uniform bound, from \rec{R08}, is far larger than the available
margin. Favorable finite rational evaluations receive much of their
saving from large numerator primes. Those factors have not been
controlled under the hypothetical equality $x=\theta$.

\subsection{Direct output rounding gives a weaker sufficient target}

There is no requirement that a successful correction have small
coefficients: the arithmetic estimate is uniform over their size.
We can therefore round directly in its $b$-dimensional output space.
This removes some unnecessary geometric requirements of the earlier
graph-lattice searches. The result combines \rec{R06} and \rec{R22}.

Write the available coefficient lattice as $U=U_0+Z\Gamma$,
where $\Gamma$ is an invertible integer $b$ by $b$ matrix. Set
$\mathcal L=\Gamma WM(\theta)^{-1}$ and $F_0=U_0\Gamma^{-1}$.
Then the relative determinant is $\det(I_b+\mathcal L(Z+F_0))$.
Choose integer input vectors $z_1,\ldots,z_b$ with independent images
$y_j=\mathcal Lz_j$, let $p_j$ be their Euclidean Gram--Schmidt
vectors, and put $d_j=\|p_j\|^2/4$.

In each affine column class $f+\Z^n$, round $-\mathcal Lf$ to the
lattice generated by the $y_j$, from the last orthogonal coordinate
to the first. Round a fractional part $t$ down with probability $1-t$
and up with probability $t$. The output error has mean zero and
covariance at most $\Sigma_0=\frac14\sum_jp_jp_j^{\mathsf T}$:
conditional means vanish and variances are at most $1/4$, so the cross
moments vanish in reverse rounding order. The input mean need not be
zero, since the determinant sees only the output.

Round the columns independently, and let $E$ be the output error
matrix. If $e_k(d)$ denotes the $k$th elementary symmetric function,
then
\begin{equation}\label{eq:rounding-moment}
 \mathbb E\det(I_b+E)=1,\qquad
 \mathbb E\det(I_b+E)^2\le\mathcal H(d):=\sum_{k=0}^b k!e_k(d).
\end{equation}
To verify the second formula, expand into principal minors. Cross
terms with different column sets vanish by independence and zero
means. For a fixed $k$-set, multilinearity and the covariance bound
give at most $k!$ times the corresponding principal minor of
$\Sigma_0$. Summing gives \eqref{eq:rounding-moment}.

Some admissible correction consequently has relative absolute value
at most $\sqrt{\mathcal H(d)}$. If all admissible corrections are
nonzero, as under rationality at the retained prime scales, this
choice is nonzero. Without that premise, there is still a nonzero
choice of size at most $\mathcal H(d)$: if $a$ is the smallest nonzero
absolute value among the finitely many outcomes $Y$, then
$a\le a\mathbb E|Y|\le\mathbb EY^2$ because $\mathbb EY=1$.

Since $\mathcal H(d)\le b!\prod_j(1+d_j)$, a sufficient uniform
condition is
\begin{equation}\label{eq:rounding-target}
 \sum_{j=1}^b\log(1+d_j)=o(n^2).
\end{equation}
It preserves the whole leading canonical margin. This condition is
weaker than bounding the largest direction and charging it $b$ times.
It has not been proved for any explicit family of directions here.
For the $23$ finite-difference bases tested at each of $n=48,96$, the
best values of $\log\mathcal H(d)/(2n^2)$ were about $.367,.389$,
respectively, against a margin about $.00001889$ on the same scale.
For rounding in a different input metric, let $v_i$ be the resulting
Gram--Schmidt input vectors and define the covariance upper bound
\[
 \Sigma=\frac14\sum_i(\mathcal L v_i)(\mathcal L v_i)^{\mathsf T}.
\]
For the same ordered input basis, the triangular-minor argument in
\rec{R22} gives $e_k(\Sigma)\ge e_k(d)$. Thus changing the metric cannot
improve this bound for the tested bases. The actual rounding covariance
also depends on the fractional parts and need not equal $\Sigma$.

\subsection{A polynomial left inverse gives an exact bound, but loses the target}\label{subsec:kernel-inverse}

Finite calculations found unexpectedly small-degree left inverses for
the pencil on the kernel of $W$. The following result explains both
what they control and why the tested way of using them cannot provide
the missing decay. This is the final reduction in \rec{R09}.

Allow the more general coefficient lattice $U=U_0+Z\Gamma$, with
$\Gamma\in\Z^{b\times b}$ invertible. Set $k=n-b$, choose rational
matrices $H_0,K_0$ with
\[
 WH_0=0,\quad WK_0=I_b,\quad T=[H_0\ K_0]\text{ invertible},
\]
and put $D(X)=M(X)H_0$, $C(X)=M(X)K_0$. The stack satisfies
\begin{equation}\label{eq:stack-kernel}
 S(X)T=
 \begin{pmatrix}D(X)&C(X)+U_0\\0&\Gamma\end{pmatrix},
 \qquad S(X)=\begin{pmatrix}M(X)+U_0W\\\Gamma W\end{pmatrix}.
\end{equation}
This displays the $k$ directions killed by the correction separately
from the $b$ directions it can change.

Suppose a polynomial matrix $E(X)$ of degree at most $d_0$ satisfies
$E(X)D(X)=I_k$. Let $\eta(E)>0$ be the rational content of all
coefficients of all $k$ by $k$ minors of $E$. If $x=a/c$ is in lowest
terms, Cauchy--Binet in $\det(E(x)D(x))=1$ gives
\[
 \frac{c^{d_0k}}{\eta(E)g_k(D(x))}\in\Z_{>0},
\]
where $g_k$ denotes rational content of the $k$-minors. Selecting all
$b$ bottom rows in \eqref{eq:stack-kernel} shows that
$|\det\Gamma|g_k(D(x))/(|\det T|\delta(x))$ is a positive integer.
Multiplying these relations proves
\begin{equation}\label{eq:inverse-content}
 \frac{C_Ec^{d_0k}}{\delta(x)}\in\Z_{>0},\qquad
 \delta(x)\le C_Ec^{d_0k},\qquad
 C_E=\frac{|\det\Gamma|}{|\det T|\eta(E)}.
\end{equation}
Thus the degree alone is not the cost; the normalization $C_E$ matters.

To apply this bound to an admissible correction class, we also require
a positive integer multiplier $a_N$ that clears every determinant polynomial
in that class:
\[
 a_NP_{N,U_0+Z\Gamma}(X)\in\Z[X]
 \qquad\text{for every }Z\in\Z^{n\times b}.
\]
Theorem~\ref{thm:arithmetic-cost} supplies this for its scalar class
$U=U_N^0+m_NZ$. Section~4 of \rec{R08} supplies it for the enlarged
admissible class, with the same leading arithmetic cost. Integrality
and invertibility of $\Gamma$ alone do not give this clearing property.

For either admissible class, there is a decisive obstruction to using
\eqref{eq:inverse-content}. At $x=0$, every cleared determinant is
integral. Finite differences then show that $a_N$ clears all mixed
minors, so $a_N\delta(0)$ is a positive integer. Applying
\eqref{eq:inverse-content} at zero gives
\begin{equation}\label{eq:bridge-obstruction}
 a_N C_E\in\Z_{>0}.
\end{equation}
Consequently, after multiplying the bound at $a/c$ by its required
clearing factor $c^na_N$, the proposed upper bound is
$(a_NC_E)c^{n+d_0k}\ge1$. It cannot certify a cleared value below one
at even a single size. Optimizing the left inverse or its coefficient
content does not change this conclusion.

The loss has an identifiable source. Selecting all bottom rows removes
$C(X)+U_0$; replacing the remaining minors by coefficient content also
removes their dependence on the particular analytic value $\theta$.
The resulting estimate holds at every rational $x$. It is a limitation
of this estimate, not a lower bound for the error of an actual correction.

\subsection{Small-degree kernel inverses do not bound correction loss}

A concrete family shows why favorable kernel data are insufficient on
their own. Let $D_e(X)$ have columns $e_j-Xe_{j+1}$, $0\le j<e$,
in $e+1$ rows. A left inverse has entries $X^{i-j}$ for $j\le i<e$
and zero otherwise, including the last column. Its degree is $e-1$,
its coefficients are zero or one, and a left null row is
$r_e(X)=(X^e,\ldots,X,1)$.

Take $6N$ chains of length six and $N$ chains of length five. This
gives precisely $n=48N$, $k=41N$, $b=7N$, with a degree-five left
inverse. For coefficient degrees $d=0,\ldots,4$, the left-multiplication
map has rank $(d+1)48N$; at $d=5$ its rank is $287N$. These are the
same profiles found at the tested moment sizes. The rank formula follows
because the kernel on one chain consists of
$f_j=X^{e-j}h$, with $\deg h\le d-e$.

Set $s=3^t$, $m=s+2$. Complete each chain by a column $m^{-1}e_e$;
the resulting square matrix has determinant $m^{-b}$. Take the
correction row on each block to be $(1/2-1/m)e_e^{\mathsf T}$ and
the initial correction column to be $e_e$. Block elimination gives
\[
 \det(M(X)+(U_0+mZ)W)=2^{-b}\det(I_b+sR_0(X)Z),
\]
where $R_0$ consists of the block rows $r_e$. At an integer $x$, the
matrices $R_0(x)Z$ range over all integer $b$ by $b$ matrices, since
each $r_e$ ends in one. The attained values are therefore
$2^{-b}(1+s\Z)$, whose least absolute value is $2^{-b}$.
The loss relative to the canonical determinant is $(m/2)^b$.
With $t=n$, its logarithm divided by $n^2$ tends to
$(7/48)\log3>0$.

At a fixed rational $x=a/c$, a clearing factor $2^bc^k$ has logarithm
$O_x(n)$. For $t>6v_3(c)$ the correction determinant is nonzero by
its reduction at $3$. Thus this example has simple kernel inverses,
small arithmetic cost and nonvanishing, yet a large correction loss.
It does not reproduce the special moment identities or their
displacement relations; it shows why the kernel profile by itself
cannot replace those missing analytic estimates.

\subsection{What successful finite approximations do and do not imply}

There is also a useful explanation for very small corrected values in
finite numerical searches. Let $M(X)=A+XB$ with rational $A,B,W$,
$\rank W=b$, and let $H_0$ span $\ker W$. Suppose
$\rank[AH_0\ BH_0]=n$ and $M(\theta)$ is invertible. Then
\begin{equation}\label{eq:conditional-density}
 \Gamma WM(\theta)^{-1}\Z^n\text{ is dense in }\R^b
 \quad\Longleftrightarrow\quad\theta\notin\Q.
\end{equation}
For rational $\theta$ the image is contained in a rational lattice.
For irrational $\theta$, a failure of density would give a nonzero
row $y$ such that $v=y\Gamma WM(\theta)^{-1}\in\Z^n$.
Then $v(AH_0+\theta BH_0)=0$. Rationality of these two matrices
forces both terms to vanish separately, and their joint rank forces
$v=0$, hence $y=0$, a contradiction. We used the elementary character
criterion: a subgroup of $\R^b$ is dense precisely when no nonzero
real linear functional takes integer values on it.

The joint-rank condition was checked at $n=48,96$ by nonzero minors
modulo $1000003$, with residues $306830,229640$. The numerical
nonsingularity at $\theta$ was checked with intervals. At these sizes,
irrationality would already imply arbitrarily fine correction with
unrestricted coefficients. This explains why a finite search can work
very well without proving the needed uniform theorem. Bounds on the
attainable determinants must still be strong enough to preserve the
decay as $N$ grows. Controlling coefficient sizes is one possible way
to obtain such bounds; Theorem~\ref{thm:arithmetic-cost} itself permits
corrections of any size.
